\section{Introduction}
Recent work has explored cost-effective methods for constructing reinforcement learning data with verifiable environments such as code, puzzles, and related tasks~\citep{zeng2025rlvescalingreinforcementlearning, wang2026rlanythingforgeenvironmentpolicy}, which yield deterministic outputs for given inputs. In this setting, models are required to reason over environment descriptions and inputs to produce correct outputs without relying on external tools~\citep{jiang-etal-2025-logicpro, ding2025longreasonarenalongreasoningbenchmark}.

Prior studies have demonstrated that increasing the number of environments leads to significant performance gains~\citep{zeng2025rlvescalingreinforcementlearning}, naturally motivating the scaling of environment pools. While several works have pursued limited scaling through automated environment synthesis~\citep{xu2026scalersyntheticscalableadaptivelearning,he2026resyn}, synthesizing environments individually increases the environment pool only linearly with construction cost. This linear scaling restricts the achievable diversity within a fixed budget and remains inadequate to support optimal reasoning generalization. Instead of generating entirely new environments from scratch, this study explores expanding the environment space by leveraging existing ones. Drawing inspiration from the closure property of transformation composition~\citep{birkhoff2017survey} and prior research on problem composition~\citep{yuan2026from,pei-etal-2025-mathfusion,chen-etal-2024-sifo}, this work introduces RACES (\textbf{R}ecursive \textbf{A}utomated \textbf{C}omposition for \textbf{E}nvironment \textbf{S}caling), a framework that composes existing environments to achieve scaling beyond linear growth.

The central insight of RACES is that when the codomain (output type) of one environment aligns with the domain (input type) of another, they can be fused into a new verifiable environment, enabling recursive composition.
As illustrated in Figure~\ref{fig:lego}, akin to LEGO bricks snapping together to form complex structures, verifiable environments can be assembled to create increasingly challenging reasoning tasks\footnote{If environments $f(x)$ and $g(x)$ each produce deterministic verifiable outputs, then their composition $g\circ f(x)=g(f(x))$ also yields a verifiable output via intermediate results.}.
To systematically convert these assembled composites into training problems, RACES defines a set of \emph{composition operators} including \textsc{SEQUENTIAL}, \textsc{PARALLEL}, \textsc{SORT} and \textsc{SELECT}, each specifying how an executable composite is presented to the model and how the response is verified.
For instance, \textsc{SEQUENTIAL} presents an assembled chain in execution order and requires the model to compute intermediate outputs, whereas \textsc{SORT} conceals the order and asks the model to recover a valid permutation that achieves a specified target output. By varying the operator choice and composition size, RACES scales to an unbounded space of environments and introduces diverse reasoning patterns, providing a robust basis for reasoning generalization.

To demonstrate the effectiveness of RACES, we construct a pool of 300 environments, implement four representative composition operators, and generate tens of thousands of composite environments for reinforcement learning on Qwen3~\citep{qwen3technicalreport} and DeepSeek-R1-Distill-Qwen~\citep{deepseekai2025deepseekr1incentivizingreasoningcapability}.
Across all model configurations, RACES consistently enhances generalization. For DeepSeek-R1-Distill-Qwen-14B, RACES increases the average score by 3.1 points (51.3 compared to 48.2), with notable improvements on IFEval (+4.0) and LongBench-v2 (+3.5). Similarly, RACES boosts Qwen3-14B performance (61.1 compared to 58.8).
Furthermore, RACES enables significantly more efficient environment utilization. Using only 50 base environments, it achieves performance comparable to training on 300 uncompounded environments.
Additional analysis demonstrates that the difficulty of composite environments can be systematically modulated by adjusting the composition size, offering greater flexibility for training design.
These findings establish RACES as an efficient and controllable approach for scaling verifiable environments and improving reasoning generalization.
